\chapter{Ideals, Prime Splitting and Class Groups}\label{ch:ideals}

\begin{center}\itshape
In $\Z[\sqrt{-5}]$, $6=2\cdot3=(1+\sqrt{-5})(1-\sqrt{-5})$: unique factorisation is lost.\\
Kummer's answer was to invent the missing ``ideal'' factors, and with them factorisation is unique again.
\end{center}

On 1 March 1847 Gabriel Lam\'e announced to the Paris Academy that he had proved Fermat's last theorem. His idea was to factor $x^p+y^p$ over the
numbers $a_0+a_1\zeta+\dots$ built from a root of unity $\zeta$. Joseph Liouville stood up and asked the awkward question: do these numbers factor uniquely
into primes? A few weeks later a letter arrived from Ernst Kummer, who had shown three years earlier that they do not, in general. Kummer's rescue was to
invent ``ideal numbers'', the missing factors, and today's ideals are their descendants. This chapter tells that story in the simplest setting, the
quadratic fields.

We use $K=\Q(\sqrt d)$ and its full ring of integers
$\mathcal O=\Z[\omega_d]$ from Chapter~\ref{ch:quadratic};
$d\ne0,1$ is squarefree. Conjugation, trace, norm and units
retain their established meanings. In particular element norms may be
negative, whereas ideal norms will be positive.

\begin{reminder}[What changes, and what does not]
The field $K$ allows division; its ring of integers $\mathcal O_K$ records which
elements behave arithmetically like integers. An element is a point; a principal
ideal is all its multiples; a general ideal is a set of allowable combinations.
Keep those three levels distinct. In particular, $(2)$ is a set, not the number
$2$, and $P^2$ is an ideal product, not a set of pointwise squares.
\end{reminder}

\section{Kummer's ideal numbers}

In $\Z[\sqrt{-5}]$ the numbers $2,3,1\pm\sqrt{-5}$ are irreducible and pairwise non-associate (Exercise~\ref{ex:zsqrt5}), and $6$ has two genuinely different factorisations. Kummer (1847) had met the same phenomenon in the rings $\Z[\zeta_p]$ while working on Fermat's equation; his idea was that $2$ and $1+\sqrt{-5}$ should have an ``ideal'' common factor, invisible among the numbers of the ring. Dedekind (1871) made this precise by replacing a number $\alpha$ by the set $(\alpha)=\alpha\mathcal O$ of its multiples, and an ``ideal number'' by a set that behaves like a set of multiples.

An \emph{ideal} of $\mathcal O$ is a subgroup $I$ of $(\mathcal O,+)$ such that $x\in I$, $y\in\mathcal O$ imply $xy\in I$. For $\alpha_1,\dots,\alpha_k\in\mathcal O$ the set $(\alpha_1,\dots,\alpha_k)=\alpha_1\mathcal O+\dots+\alpha_k\mathcal O$ is an ideal, the ideal \emph{generated} by the $\alpha_i$; an ideal $(\alpha)$ is \emph{principal}. Note that $(\alpha)\subseteq(\beta)$ means $\beta\mid\alpha$, and $(\alpha)=(\beta)$ means $\alpha,\beta$ are associates: principal ideals are numbers up to units. The \emph{product} $IJ$ of two ideals is the set of finite sums $\sum x_iy_i$ with $x_i\in I$, $y_i\in J$. It is an ideal, the product is commutative and associative with neutral element $\mathcal O=(1)$, and expanding products of linear combinations shows
\[
\begin{gathered}(\alpha_1,\dots,\alpha_k)(\beta_1,\dots,\beta_l)=(\alpha_i\beta_j)_{i,j},\qq\text{in particular }(\alpha)(\beta)=(\alpha\beta).\end{gathered}
\]
Finally $\bar I=\{\bar x:x\in I\}$ is an ideal (as $\sigma$ preserves $\mathcal O$), with $\overline{(\alpha,\beta)}=(\bar\alpha,\bar\beta)$ and $\overline{IJ}=\bar I\bar J$.

\begin{think}
The only common divisors of $2$ and $1+\sqrt{-5}$ in $\Z[\sqrt{-5}]$ are $\pm1$ (compare norms). Yet the set $P=\{2x+(1+\sqrt{-5})y\}$ of their combinations is not all of $\Z[\sqrt{-5}]$. Which points $a+b\sqrt{-5}$ does it contain? Draw it before looking at Figure~\ref{fig:lattice}.
\end{think}

\begin{lemma}\label{lem:twogen}
Every ideal $I\ne\{0\}$ is of the form $\Z\alpha+\Z\beta$, with $\alpha,\beta$ linearly independent over $\Q$; in particular $I=(\alpha,\beta)$ is generated by two elements.
\end{lemma}
\begin{proof}
Identify $\mathcal O$ with $\Z^2$ via $m+n\omega\mapsto(m,n)$ (Proposition~\ref{prop:OK}). Choose $0\ne\eta\in I$. Then $N\eta=\eta\bar\eta\in I\cap\Z$ is non-zero, and $(N\eta)\omega\in I$ has non-zero second coordinate. The second coordinates of elements of $I$ form a non-zero subgroup of $\Z$, so they are the multiples of some $c>0$; pick $\beta\in I$ with second coordinate $c$. The elements of $I$ with second coordinate $0$ form a subgroup $a\Z$ of $\Z$; let $\alpha=a>0$ (the nonzero integer $N\eta$ lies in this subgroup). Since $\beta$ has nonzero second coordinate, $\alpha,\beta$ are linearly independent. Any $x\in I$ has second coordinate $kc$, and $x-k\beta\in I$ has second coordinate $0$, so it is a multiple of $\alpha$. Hence $I=\Z\alpha+\Z\beta$. Then $(\alpha,\beta)\subseteq I$ since $\alpha,\beta\in I$, and $I=\Z\alpha+\Z\beta\subseteq(\alpha,\beta)$.
\end{proof}

\begin{proposition}[Kummer's ideal factors of $6$]\label{prop:six}
In $\mathcal O=\Z[\sqrt{-5}]$, let $P=(2,1+\sqrt{-5})$, $Q=(3,1+\sqrt{-5})$ and $\bar Q=(3,1-\sqrt{-5})$. Then
\[
\begin{gathered}P^2=(2),\qquad Q\bar Q=(3),\qq PQ=(1+\sqrt{-5}),\qquad P\bar Q=(1-\sqrt{-5}),\end{gathered}
\]
so both factorisations of $6$ are regroupings of $(6)=P\cdot P\cdot Q\cdot\bar Q$. The ideals $P$, $Q$, $\bar Q$ are not principal.
\end{proposition}
\begin{proof}
First $\bar P=(2,1-\sqrt{-5})=P$, since $1-\sqrt{-5}=2-(1+\sqrt{-5})$. Then
$P^2=(4,\,2+2\sqrt{-5},\,(1+\sqrt{-5})^2)=(4,\,2+2\sqrt{-5},\,-4+2\sqrt{-5})$; every generator is divisible by $2$, and $2=(2+2\sqrt{-5})-(-4+2\sqrt{-5})-4$ lies in $P^2$, so $P^2=(2)$.
$Q\bar Q=(9,\,3-3\sqrt{-5},\,3+3\sqrt{-5},\,6)$: all generators are multiples of $3$, and $3=9-6$, so $Q\bar Q=(3)$.
$PQ=(6,\,2+2\sqrt{-5},\,3+3\sqrt{-5},\,-4+2\sqrt{-5})$: each generator is a multiple of $1+\sqrt{-5}$ (as $6=(1+\sqrt{-5})(1-\sqrt{-5})$ and $-4+2\sqrt{-5}=(1+\sqrt{-5})^2$), and $1+\sqrt{-5}=(3+3\sqrt{-5})-(2+2\sqrt{-5})$, so $PQ=(1+\sqrt{-5})$. Conjugating, $P\bar Q=\bar P\bar Q=\overline{PQ}=(1-\sqrt{-5})$.

Every element $a+b\sqrt{-5}$ of $P$ has $a\equiv b\pmod2$: this holds for $2x$ and for $(1+\sqrt{-5})(x+y\sqrt{-5})=(x-5y)+(x+y)\sqrt{-5}$, and is preserved by sums. So $1\notin P$. If $P=(\alpha)$, then $\alpha$ divides $2$ and $1+\sqrt{-5}$, so $N\alpha$ divides $4$ and $6$; $N\alpha=a^2+5b^2=2$ is impossible, so $N\alpha=1$, $\alpha=\pm1$ and $P=\mathcal O$, a contradiction. Likewise, if $Q=(\alpha)$ then $N\alpha\mid\gcd(9,6)=3$; $a^2+5b^2=3$ is impossible and $Q=\mathcal O$ would give $1=3x+(1+\sqrt{-5})y$, whence, multiplying by $1-\sqrt{-5}$, $3\mid1-\sqrt{-5}$, which is false. The same holds for $\bar Q$ by conjugation.
\end{proof}

\begin{figure}[ht]
\centering
\fig{ch11-lattice}
\caption{$\Z[\sqrt{-5}]$ is a lattice in $\C$, and an ideal is a sublattice stable under multiplication by $\sqrt{-5}$. The ideal $P$ (blue) contains every second point, the ideal $(2)=P^2$ (circled) every fourth. $P$ is a perfectly good ``number'', it just is not the set of multiples of any one element.}
\label{fig:lattice}
\end{figure}

\begin{remark}
In Figure~\ref{fig:lattice}, the proportions $\frac12$ and $\frac14$ are $\frac1{N(P)}$ and $\frac1{N((2))}$, with the norm defined in the next section. In general $N(I)$ is the index of $I$ in $\mathcal O$, the number of classes of $\mathcal O$ modulo $I$; we prove this identity in Proposition~\ref{prop:t1-norm-index} below, after establishing ideal factorization.
\end{remark}

Two elements are congruent modulo $I$ when their difference belongs
to $I$. These equivalence classes form the \emph{quotient ring}
$\mathcal O/I$: addition and multiplication of representatives are
well defined because $I$ absorbs multiplication by every ring element.
Indeed, replacing $a,b$ by $a+i,b+j$ changes their sum by $i+j$
and their product by $aj+bi+ij$, all in $I$.
The quotient need not be a field. For $I=(2)$ in $\Z[\sqrt{-5}]$,
the nonzero class of $1+\sqrt{-5}$ has square zero.

\subsection*{Seeing an ideal through its remainders}
Set $s=\sqrt{-5}$ and $P=(2,1+s)$. The simple rule
\[
\rho:\Z[s]\longrightarrow\Z/2\Z,
\qquad \rho(a+bs)=a+b\pmod2
\]
preserves sums and products: modulo $2$, the relation $s^2=-5$ becomes
$1^2=1$. Its kernel is exactly $P$. One inclusion follows from the generators;
for the other, if $a+b$ is even, then
\[
a+bs=\frac{a-b}{2}\,2+b(1+s)\in P.
\]
Thus there are exactly two remainders modulo $P$, represented by $0,1$.
Modulo $(2)$, however, the coefficients $a,b$ vary independently, giving four
remainders $0,1,s,1+s$.

\begin{figure}[ht]
\centering\fig{ch11-quotients}
\caption{Four coefficient remainders modulo $(2)$ merge into two remainders
modulo $P=(2,1+\sqrt{-5})$. Adding $1+\sqrt{-5}$ pairs the two solid dots and
pairs the two open dots. Passing to the larger ideal identifies more points.}
\label{fig:ideal-quotients}
\end{figure}

This finite arithmetic also explains why $P$ is maximal: if an ideal contains
$P$ and any element outside it, that element has remainder $1$, so the larger
ideal contains $1$. The norm formalism below will generalise this observation.

\begin{think}
Why is ``$P$ is not principal'' different from ``$P$ cannot be generated''?
It \emph{is} generated, by $2$ and $1+s$. The point is that no single element
can generate it. Ideals have not broken arithmetic; they have relaxed its
one-person staffing policy.
\end{think}

\portrait{minkowski}{Hermann Minkowski (1864--1909). Photograph from \emph{Raum und Zeit}, 1909. Public domain, via Wikimedia Commons.}

\section{Unique factorisation of ideals}

Everything rests on one identity, a version of Gauss' lemma for ideals; the idea goes back to Dedekind and Hurwitz.

\begin{lemma}\label{lem:hurwitz}
For $\alpha,\beta\in\mathcal O$ not both zero, $(\alpha,\beta)(\bar\alpha,\bar\beta)=(n)$, where $n\ge1$ is the gcd in $\Z$ of $N\alpha$, $N\beta$ and $\operatorname{Tr}(\alpha\bar\beta)$.
\end{lemma}
\begin{proof}
The product is generated by $\alpha\bar\alpha=N\alpha$, $\alpha\bar\beta$, $\bar\alpha\beta$ and $\beta\bar\beta=N\beta$. It contains $N\alpha$, $N\beta$ and $\alpha\bar\beta+\bar\alpha\beta=\operatorname{Tr}(\alpha\bar\beta)$, hence their gcd $n$ (a $\Z$-combination of them by B\'ezout); $n\ne0$ since $N\alpha$ or $N\beta$ is non-zero. Conversely $n$ divides each generator in $\mathcal O$: clearly $N\alpha$ and $N\beta$; and $\gamma=\frac{\alpha\bar\beta}n$ has $\operatorname{Tr}\gamma=\frac{\operatorname{Tr}(\alpha\bar\beta)}n\in\Z$ and $N\gamma=\frac{N\alpha}n\cdot\frac{N\beta}n\in\Z$, so $\gamma\in\mathcal O$ by Proposition~\ref{prop:intcrit}, and $\bar\gamma=\frac{\bar\alpha\beta}n\in\mathcal O$ too. So the product is $(n)$.
\end{proof}

By Lemma~\ref{lem:twogen}, every ideal $I\ne\{0\}$ satisfies $I\bar I=(n)$ for some integer $n\ge1$, and $n$ is unique: $(n)=(m)$ with $n,m\ge1$ gives $\frac nm,\frac mn\in\mathcal O\cap\Q=\Z$, so $n=m$. This $n$ is the \emph{norm} $N(I)$ of $I$. From now on ``ideal'' means non-zero ideal.

\begin{proposition}\label{prop:idealrules}
Let $I,J,L$ be ideals and $m\ge1$ an integer.
\begin{enumerate}[label=(\alph*)]
\item $N(IJ)=N(I)N(J)$, $N((\alpha))=|N\alpha|$, $N(I)\in I$, and $N(I)=1$ if and only if $I=\mathcal O$.
\item (Cancellation) If $IJ=IL$ then $J=L$.
\item (To contain is to divide) $I\subseteq J$ if and only if $I=JL$ for some ideal $L$; then $N(J)\mid N(I)$.
\item Only finitely many ideals contain $m$.
\end{enumerate}
\end{proposition}
\begin{proof}
(a) $(IJ)\overline{(IJ)}=I\bar I\,J\bar J=(N(I))(N(J))=(N(I)N(J))$. Also $(\alpha)(\bar\alpha)=(N\alpha)=(|N\alpha|)$. Next, $N(I)\in I\bar I\subseteq I$ because $\bar I\subseteq\mathcal O$. If $N(I)=1$ then $1\in I$ and $I=\mathcal O$; conversely $N(\mathcal O)=1$ since $\mathcal O\cdot\mathcal O=(1)$.

(b) Multiply by $\bar I$: $N(I)J=N(I)L$ as sets, and dividing every element by $N(I)$ gives $J=L$.

(c) If $I=JL$ then $I\subseteq J$, since $J$ absorbs multiplication. Conversely let $I\subseteq J$ and $k=N(J)$. The ideal $I\bar J$ is contained in $J\bar J=k\mathcal O$, so $L=\frac1k\,I\bar J$ is contained in $\mathcal O$ and is an ideal; and $JL=\frac1k\,(J\bar J)I=I$. Then $N(I)=N(J)N(L)$ by (a).

(d) An ideal $I\ni m$ contains $m\mathcal O$, so it is a union of classes of $\mathcal O/m\mathcal O$, which has $m^2$ elements ($a+b\omega$ modulo $m$ is determined by $a,b$ modulo $m$). There are at most $2^{m^2}$ such unions.
\end{proof}

An ideal $P$ is \emph{maximal} if $P\ne\mathcal O$ and the only ideals containing $P$ are $P$ and $\mathcal O$. These are the irreducible ``ideal numbers''.

\begin{lemma}\label{lem:maximal}
\begin{enumerate}[label=(\alph*)]
\item If $P$ is maximal and $P\supseteq IJ$, then $P\supseteq I$ or $P\supseteq J$. In particular, if $xy\in P$ then $x\in P$ or $y\in P$.
\item Every ideal $I\ne\mathcal O$ is contained in a maximal ideal.
\item An ideal whose norm is a prime number is maximal.
\end{enumerate}
\end{lemma}
\begin{proof}
(a) Let $a\in I\setminus P$. The ideal $P+a\mathcal O$ strictly contains $P$, so it is $\mathcal O$: $1=p+ax$ with $p\in P$. For $b\in J$, $b=pb+(ab)x\in P$, because $ab\in IJ\subseteq P$. So $J\subseteq P$. Apply it to $I=(x)$, $J=(y)$ for the second claim.

(b) The ideals containing $I$ contain $N(I)$, so they are finitely many by Proposition~\ref{prop:idealrules}(d). Among those different from $\mathcal O$ (there is at least $I$), choose one that is not strictly contained in another: it is maximal.

(c) If $P\subsetneq J\subsetneq\mathcal O$, write $P=JL$; then $L\ne\mathcal O$ (else $P=J$), so $N(P)=N(J)N(L)$ with both factors $\ge2$.
\end{proof}

\begin{theorem}[Dedekind]\label{thm:dedekind}
Every ideal $I\ne\mathcal O$ of $\mathcal O$ is a product of maximal ideals, uniquely up to order.
\end{theorem}
\begin{proof}
Existence, by strong induction on $N(I)\ge2$. Take a maximal $P\supseteq I$ and write $I=PL$. Then $N(L)=\frac{N(I)}{N(P)}<N(I)$ since $N(P)\ge2$; either $L=\mathcal O$ and $I=P$, or $L$ is a product of maximal ideals by induction.

Uniqueness: let $P_1\cdots P_r=Q_1\cdots Q_s$ with all factors maximal. By Lemma~\ref{lem:maximal}(a), applied repeatedly, $P_1\supseteq Q_1\cdots Q_s$ gives $P_1\supseteq Q_j$ for some $j$, and then $P_1=Q_j$ because $Q_j$ is maximal and $P_1\ne\mathcal O$. Cancel it (Proposition~\ref{prop:idealrules}(b)) and induct on $r$. When $r=0$ the left side is $\mathcal O$, which cannot equal a product with $s\ge1$ factors, since such a product lies in $Q_1\ne\mathcal O$.
\end{proof}

So Kummer was right: with ideals, arithmetic in $\mathcal O$ is as good as in $\Z$. As in $\Z$, each ideal has a valuation $v_P(I)$ at each maximal $P$, and if $I$ and $J$ have no maximal factor in common and $IJ$ is a $k$-th power, then $I$ and $J$ are $k$-th powers (compare the exponents of each $P$).

\section{How a prime number splits}

\begin{think}
Figure~\ref{fig:split} colours the primes $p<240$ by the way the ideal $(p)$ factors in $\Z[\sqrt{-5}]$. Before reading the theorem, guess the rule: does the colour depend only on $p$ modulo some number? Which one? (A green prime is $x^2+5y^2$, for instance $29=3^2+5\cdot2^2$.)
\end{think}

\begin{figure}[ht]
\centering
\fig{ch11-split}
\caption{How the primes $p<240$ factor in $\Z[\sqrt{-5}]$, computed with Theorem~\ref{thm:split} and Exercise~\ref{ex:fermat5}.}
\label{fig:split}
\end{figure}

\begin{theorem}\label{thm:split}
Let $p$ be an odd prime and $\left(\frac dp\right)$ the Legendre symbol.
\begin{enumerate}[label=(\alph*)]
\item (split) If $\left(\frac dp\right)=1$ and $a^2\equiv d\pmod p$, then $(p)=P\bar P$ with $P=(p,a+\sqrt d)$, $N(P)=N(\bar P)=p$ and $P\ne\bar P$.
\item (ramified) If $p\mid d$, then $(p)=P^2$ with $P=(p,\sqrt d)$ and $N(P)=p$.
\item (inert) If $\left(\frac dp\right)=-1$, then $(p)$ is maximal.
\end{enumerate}
Every maximal ideal contains exactly one prime number $p$, so it occurs in the factorisation of exactly one $(p)$.
\end{theorem}
\begin{proof}
(a) By Lemma~\ref{lem:hurwitz}, $P\bar P=(n)$ with $n=\gcd\big(p^2,\,a^2-d,\,\operatorname{Tr}(p(a-\sqrt d))\big)=\gcd(p^2,a^2-d,2ap)$. Now $p$ divides all three, and $p^2\nmid2ap$ because $p\nmid2a$ ($p\nmid d$ gives $p\nmid a$). So $n=p$, and $N(P)=N(\bar P)=p$ (conjugate ideals have the same norm, as $\bar P\,\overline{\bar P}=P\bar P$). If $P=\bar P$, it would contain $(a+\sqrt d)+(a-\sqrt d)=2a$ and $p$, which are coprime integers, so $P=\mathcal O$ and $N(P)=1$: impossible.

(b) Here $n=\gcd(p^2,N\sqrt d,\operatorname{Tr}(p\cdot(-\sqrt d)))=\gcd(p^2,-d,0)=p$ since $d$ is squarefree. As $\bar P=(p,-\sqrt d)=P$, this reads $P^2=(p)$.

(c) Let $J\supseteq(p)$ with $J\ne\mathcal O$, and write $(p)=JL$, so $N(J)N(L)=p^2$ and $N(J)\in\{p,p^2\}$. If $N(J)=p^2$ then $N(L)=1$, $L=\mathcal O$ and $J=(p)$. Suppose $N(J)=p$; then $J\ne(p)$, so pick $\alpha\in J\setminus p\mathcal O$. Then $N\alpha=\alpha\bar\alpha\in J\bar J=(p)$, and $N\alpha\in p\mathcal O\cap\Z=p\Z$ (Proposition~\ref{prop:intcrit}). Write $\alpha=\frac{u+v\sqrt d}2$ with $u,v\in\Z$: $p\mid u^2-dv^2$. If $p\mid v$, then $p\mid u$, and $\frac\alpha p=\frac{u/p+(v/p)\sqrt d}2$ satisfies the parity conditions of Proposition~\ref{prop:OK} (since $p$ is odd, $\frac up\equiv u$ and $\frac vp\equiv v\pmod2$), so $\alpha\in p\mathcal O$: excluded. So $p\nmid v$ and $d\equiv(uv^{-1})^2\pmod p$, contradicting $\left(\frac dp\right)=-1$.

Finally, a maximal $P$ contains $N(P)\ge2$, hence (Lemma~\ref{lem:maximal}(a)) one of its prime factors $p$. If it contained two distinct primes $p,q$, it would contain $1=up+vq$. And $P\ni p$ means $P\mid(p)$ by Proposition~\ref{prop:idealrules}(c).
\end{proof}

For $p=2$ the answer depends on $d\bmod8$ (Exercise~\ref{ex:two}). In $\Z[i]$, $\left(\frac{-1}p\right)=1$ if and only if $p\equiv1\pmod4$, and Theorem~\ref{thm:split} recovers Theorem~\ref{thm:gaussprimes} for odd $p$. The general principle: \emph{the factorisation of $p$ is decided by a Legendre symbol}, hence, by quadratic reciprocity, by $p$ modulo $4|d|$.

\subsection*{Three behaviours, checked on actual primes}
In $\Z[\sqrt{-5}]$, the polynomial $X^2+5$ provides a quick preview of the
three cases of Theorem~\ref{thm:split}.
\begin{itemize}
\item Modulo $3$, $X^2+5=X^2-1$ has the two distinct roots $1,-1$.
The two ideals $(3,1+s)$ and $(3,1-s)$ are distinct and their product is $(3)$.
Neither is principal: an element generating one would have norm $3$, whereas
$a^2+5b^2=3$ has no solution.
\item Modulo $5$, the polynomial is $X^2$, with a repeated root. The ideal
$(5,s)=(s)$ occurs twice in $(5)=(s)^2$; the sign $s^2=-5$ does not affect
principal ideals.
\item Modulo $11$, $-5\equiv6$ is not a square. The ideal $(11)$ stays maximal,
and its norm is $11^2$, not $11$.
\end{itemize}
Inert primes are a useful warning: a maximal ideal need not have prime norm.
Its norm can be a prime \emph{power}. The lemma only asserted the implication
``prime norm implies maximal'', not its converse.

\begin{reminder}[The containment sign points the other way]
$P\mid I$ means $I=PL$ for an integral ideal $L$, hence $I\subseteq P$.
For ordinary numbers the model is $(6)\subseteq(2)$ because $2\mid6$.
When checking an ideal argument, quietly replace every containment by that
example. It saves a remarkable amount of mental gymnastics.
\end{reminder}

\section{Ideal arithmetic and the number of remainders}

Ideal factorization now supplies a useful dictionary. For each maximal $P$,
\[
v_P(I+J)=\min(v_P(I),v_P(J)),\qquad
v_P(I\cap J)=\max(v_P(I),v_P(J)).
\]
Here $I+J=\{a+b:a\in I,b\in J\}$.
To prove the first identity, let $G$ be the ideal with the minimum exponents.
Both $I,J$ lie in $G$, so $I+J\subseteq G$. Conversely any ideal containing
both $I,J$ has exponents at most those minima, and therefore contains $G$.
Apply this to $I+J$ itself. The intersection identity follows by reversing
the containments: it is the largest ideal contained in both.
All these steps use the containment/divisibility rule already proved.

For distinct maximal ideals $P,Q$, this gives $P+Q=\mathcal O$.
More generally, when $I+J=\mathcal O$, choose $u\in I,v\in J$
with $u+v=1$. The map
\[
\mathcal O/(IJ)\longrightarrow\mathcal O/I\times\mathcal O/J
\]
given by reducing both coordinates is an isomorphism:
$I\cap J=IJ$ gives injectivity, while the element $av+bu$
has residues $a$ modulo $I$ and $b$ modulo $J$, giving surjectivity.
This proves the Chinese remainder theorem for these ideals from the
same switches as in Chapter~\ref{ch:congruences}.

\begin{proposition}[Norm counts remainders]\label{prop:t1-norm-index}
For every nonzero ideal $I$, the number of classes in $\mathcal O/I$
is $N(I)$.
\end{proposition}
\begin{proof}
First let $P$ be maximal. It contains exactly one rational prime $p$,
so $(p)=PL$ and $N(P)N(L)=p^2$. Thus $N(P)$ is $p$ or $p^2$.
If it is $p^2$, then $L=\mathcal O$ and $P=(p)$, whose quotient has
$p^2$ coefficient remainders. If it is $p$, then
$(p)\subsetneq P\subsetneq\mathcal O$. The quotient
$\mathcal O/(p)$ is a two-dimensional vector space over $\Z/p\Z$,
and $P/(p)$ is a nonzero proper subspace, hence one-dimensional.
There are therefore $p$ cosets of it, giving
$\#(\mathcal O/P)=p=N(P)$.

Fix $k\ge0$ and choose $a\in P^k\setminus P^{k+1}$, possible because
their norms differ. The sum rule above gives
$(a)+P^{k+1}=P^k$: the exponent of $P$ is $k$, and all other
minimum exponents are zero. Thus the additive map
$x\mapsto ax$ from $\mathcal O$ onto $P^k/P^{k+1}$
is surjective. Its kernel is $P$, since the exponent of $P$
in $(a)$ is exactly $k$. Consequently
$\#(P^k/P^{k+1})=N(P)$.
Counting cosets successively in
$\mathcal O\supset P\supset\cdots\supset P^e$ gives
$\#(\mathcal O/P^e)=N(P)^e$.
Finally factor $I=\prod P_i^{e_i}$ and use the just-proved
Chinese remainder theorem for the pairwise coprime prime powers.
The quotient size is $\prod N(P_i)^{e_i}=N(I)$.
\end{proof}

\subsection*{A product records sums of products}
Return to $P=(2,1+s)$ with $s=\sqrt{-5}$.
Its product with itself contains every finite sum of products,
so it is generated by $4,2+2s,-4+2s$.
The combination
$(2+2s)-(-4+2s)-4=2$ proves $(2)\subseteq P^2$,
and divisibility of all three generators by $2$ proves the reverse.
An elementwise square such as $(1+s)^2$ alone would not prove this identity.

Likewise $Q=(3,1+s)$ gives $PQ=(1+s)$, of norm $6$.
Their sum is $\mathcal O$, since $2\in P$ and $3\in Q$;
hence $P\cap Q=PQ$ and a remainder modulo $1+s$ is exactly
a pair of remainders modulo $P$ and $Q$, with $2\cdot3=6$ choices.
For example $3$ is $1$ modulo $P$ and $0$ modulo $Q$;
$-2$ reverses those roles. The same two-switch calculation has
survived from ordinary congruences to a ring without element factorization.

\section{The class group}

Two ideals are \emph{equivalent}, $I\sim J$, if $\alpha I=\beta J$ for some non-zero $\alpha,\beta\in\mathcal O$. This is an equivalence relation (for transitivity, $\alpha I=\beta J$ and $\gamma J=\delta L$ give $\alpha\gamma I=\beta\delta L$), compatible with products. An ideal is equivalent to $\mathcal O$ if and only if it is principal: if $\alpha I=(\beta)$, then $\beta\in\alpha I$, so $\frac\beta\alpha\in I\subseteq\mathcal O$ and $I=\left(\frac\beta\alpha\right)$. The classes therefore form a commutative group $\operatorname{Cl}(K)$, the \emph{class group}: the neutral element is the class of principal ideals, and the inverse of the class of $I$ is that of $\bar I$, since $I\bar I=(N(I))$. Its size $h_K$ is the \emph{class number}; it measures how far $\mathcal O$ is from being a ring where every ideal is a number.

\begin{lemma}\label{lem:approx}
There is an integer $M\ge1$, depending on $d$, such that for every $\gamma\in K$ there are $t\in\{1,\dots,M\}$ and $\kappa\in\mathcal O$ with $|N(t\gamma-\kappa)|<1$.
\end{lemma}
\begin{proof}
Let $L$ be an integer with $L^2>1+|d|$ and $M=L^2$. Write $\gamma=s+s'\sqrt d$ with $s,s'\in\Q$ and look at the $M+1$ points $(\{ks\},\{ks'\})\in[0,1)^2$, $0\le k\le M$, where $\{\cdot\}$ is the fractional part. Cut $[0,1)^2$ into $L^2=M$ squares of side $\frac1L$: two points $k_1<k_2$ fall in the same square (pigeonhole). With $t=k_2-k_1\in\{1,\dots,M\}$, $ts=x+\sigma$ and $ts'=y+\tau$ with $x,y\in\Z$ and $|\sigma|,|\tau|<\frac1L$. Then $\kappa=x+y\sqrt d\in\mathcal O$ and $|N(t\gamma-\kappa)|=|\sigma^2-d\tau^2|\le\sigma^2+|d|\tau^2<\frac{1+|d|}{L^2}<1$.
\end{proof}

\begin{theorem}\label{thm:classfinite}
With $M$ as in Lemma~\ref{lem:approx}, every ideal is equivalent to an ideal containing $M!$. In particular the class number $h_K$ is finite.
\end{theorem}
\begin{proof}
Let $I$ be an ideal and $\beta\in I$ non-zero with $|N\beta|$ least (a positive integer). For $\alpha\in I$, apply the lemma to $\gamma=\frac\alpha\beta$: $|N(t\alpha-\kappa\beta)|=|N\beta|\cdot|N(t\gamma-\kappa)|<|N\beta|$, and $t\alpha-\kappa\beta\in I$, so $t\alpha=\kappa\beta$ by minimality. As $t\mid M!$, every $\alpha\in I$ has $M!\,\alpha\in(\beta)$. Hence $J=\frac{M!}\beta I$ is contained in $\mathcal O$; it is an ideal, equivalent to $I$ ($\beta J=M!\,I$), and it contains $M!$ (take $\alpha=\beta$). By Proposition~\ref{prop:idealrules}(d) only finitely many ideals contain $M!$, so there are finitely many classes.
\end{proof}

\begin{proposition}\label{prop:h1}
If $h_K=1$, then $\mathcal O$ has unique factorisation: every non-zero non-unit is a product of irreducibles, uniquely up to order and units.
\end{proposition}
\begin{proof}
Existence holds in any of our rings, by strong induction on $|N\alpha|$, exactly as in Theorem~\ref{thm:ufd}(c). If $\pi$ is irreducible, $(\pi)$ is maximal: if $(\pi)\subseteq J\ne\mathcal O$, then $J=(\delta)$ since $h_K=1$, and $\pi=\delta x$ with $\delta$ not a unit, so $x$ is a unit and $J=(\pi)$. Now if $\pi_1\cdots\pi_r=u\rho_1\cdots\rho_s$ are two factorisations, then $(\pi_1)\cdots(\pi_r)=(\rho_1)\cdots(\rho_s)$ are two factorisations into maximal ideals; by Theorem~\ref{thm:dedekind}, $r=s$ and after reordering $(\pi_i)=(\rho_i)$, that is, $\pi_i$ and $\rho_i$ are associates.
\end{proof}

The converse also holds (Exercise~\ref{ex:ufdh1}): $h_K$ is exactly the obstruction to unique factorisation. For $d=-1,-2,-3$, Chapter~\ref{ch:gaussian} gave a Euclidean division, so all ideals are principal (the proof of Theorem~\ref{thm:ufd}(a) shows it) and $h_K=1$. For $d=-5$ we can now compute everything.

\begin{think}
In $\Z[\sqrt{-5}]$ we know one non-principal ideal, $P=(2,1+\sqrt{-5})$, and another one, $Q=(3,1+\sqrt{-5})$. Is $Q$ equivalent to $P$? (Look at $PQ$ in Proposition~\ref{prop:six}.) How many classes would you guess there are?
\end{think}

\begin{theorem}\label{thm:hm5}
The class number of $\Q(\sqrt{-5})$ is $2$: every ideal of $\Z[\sqrt{-5}]$ is either principal or equivalent to $P=(2,1+\sqrt{-5})$.
\end{theorem}
\begin{proof}
We sharpen Lemma~\ref{lem:approx}: \emph{for every $\gamma\in K$ there are $t\in\{1,2\}$ and $\kappa\in\mathcal O$ with $N(t\gamma-\kappa)<1$.} Write $\gamma=s+s'\sqrt{-5}$ and round: $s=x+\sigma$, $s'=y+\tau$ with $x,y\in\Z$ and $|\sigma|,|\tau|\le\frac12$. If $\tau^2<\frac3{20}$, then $N(\gamma-x-y\sqrt{-5})=\sigma^2+5\tau^2<\frac14+\frac34=1$. Otherwise $|\tau|\ge\sqrt{3/20}>3/8$, while $|\tau|\le1/2$. Hence $2\tau$ lies at distance strictly less than $1/4$ from $1$ or $-1$. Rounding the real coefficient of $2\gamma$ to within $1/2$ and its $\sqrt{-5}$ coefficient to that nearby integer gives $N<1/4+5/16=9/16<1$.

Now let $I$ be an ideal and $\beta\in I$ non-zero of least norm. As in the proof of Theorem~\ref{thm:classfinite}, $2\alpha\in(\beta)$ for all $\alpha\in I$, so $J=\frac2\beta I$ is an ideal equivalent to $I$ that contains $2$. By Proposition~\ref{prop:idealrules}(c), $J$ divides $(2)=P^2$, so by Theorem~\ref{thm:dedekind} $J\in\{\mathcal O,P,P^2\}$. Since $P^2=(2)$ is principal, $I$ is principal or equivalent to $P$; and $P$ is not principal (Proposition~\ref{prop:six}). So there are exactly two classes.
\end{proof}

\begin{remark}
Gauss conjectured, and Heegner (1952), Baker and Stark (1966--67) proved, that the imaginary quadratic fields with class number $1$ are exactly those with $d=-1,-2,-3,-7,-11,-19,-43,-67,-163$. The last one explains why $e^{\pi\sqrt{163}}$ is so close to an integer:
\[\begin{aligned}e^{\pi\sqrt{163}}&=262537412640768743.\eqbr\quad99999999999925\ldots\end{aligned}\]
 Kummer proved Fermat's Last Theorem for every prime $p$ that does not divide the class number of $\Q(\zeta_p)$ (the \emph{regular} primes, such as all $p<37$), by the method of Theorem~\ref{thm:flt3} with ideals in place of numbers. Both results lie beyond this series; the analytic class number formula, which computes $h_K$ with an $L$-function, is in Volume~12.
\end{remark}

\subsection*{The class group as a parity ledger}
For $K=\Q(\sqrt{-5})$, write the two ideal classes as $1,c$ with $c^2=1$.
An ideal product is principal exactly when it contains an even number of
non-principal factors, counted with multiplicity. For the factors of $6$,
\[
[P]=[Q]=[\bar Q]=c,\qquad
[P^2]=[PQ]=[P\bar Q]=[Q\bar Q]=1.
\]
The two element factorisations of $6$ group these same ideal factors into
principal pairs in different ways. The class group remembers the obstruction
that the norm alone cannot see.

A particularly useful deduction is this: if $I^3$ is principal, then $I$ is
principal here, because $[I]^2=[I]^3=1$ gives $[I]=1$.
It would be false to make that deduction in a class group with an element of
order $3$. Exercise~\ref{ex:x2p5} uses precisely this odd-exponent advantage.

\begin{think}
If $I^2$ is principal, must $I$ be principal? The ideal $P$ answers ``no''.
If both $I^2$ and $I^3$ are principal, must $I$ be principal in \emph{any}
class group? Yes: $[I]=[I]^3([I]^2)^{-1}=1$.
\end{think}

\section{A brief bridge to binary quadratic forms}

A binary quadratic form is $f(u,v)=Au^2+Buv+Cv^2$ with integer
coefficients. Its discriminant is $D=B^2-4AC$.
If $A>0$ and $D<0$, completing the square gives
\[
f(u,v)=A\left(u+\frac{Bv}{2A}\right)^2
       +\frac{-D}{4A}v^2,
\]
so it is positive except at $(0,0)$.
Here are two forms of discriminant $-20$:
\[
u^2+5v^2,\qquad 2u^2+2uv+3v^2.
\]
The first is the element norm in $\Z[\sqrt{-5}]$.
The second comes from the nonprincipal ideal $P=(2,1+s)$:
\[
N\bigl(2u+(1+s)v\bigr)
=4u^2+4uv+6v^2=2(2u^2+2uv+3v^2).
\]
Both the factor $2=N(P)$ and the middle coefficient have an arithmetic reason.

For an odd split prime $p\ne5$, let $J$ be one ideal factor above $p$.
If $J$ is principal, $p$ is represented by the first form.
If $J$ is nonprincipal, $PJ$ is principal because the class group has
order $2$: write $PJ=(\alpha)$. Now $\alpha\in P$, so
$\alpha=2u+(1+s)v$, and $N\alpha=2p$ proves
$p=2u^2+2uv+3v^2$.

Conversely, if the second form represents such a prime, the corresponding
$\alpha\in P$ has norm $2p$. Write $(\alpha)=PJ$ using the containment
rule. Then $N(J)=p$, and $[J]=[P]^{-1}=[P]$ is nonprincipal.
Thus the two forms distinguish the two ideal classes above split primes.
For example $3=f(0,1)$ and $7=f(1,1)$ belong to the second form,
whereas $29=3^2+5\cdot2^2$ belongs to the first.
A split prime cannot occur in both forms: conjugation makes its two
ideal factors either both principal or both nonprincipal.

This is one concrete instance of a wider correspondence.
We have not defined equivalence, reduction or composition of general
quadratic forms, and no such theory is used in these arguments.
Cox's book in the sources develops that next step.

\section{Exercises}

\begin{exercise}[$\star\star$]\label{ex:sqrt10}
In $\Z[\sqrt{10}]$, show that $2$, $5$ and $\sqrt{10}$ are irreducible, and deduce from $10=2\cdot5=\sqrt{10}^2$ that $h_{\Q(\sqrt{10})}\ge2$.
\end{exercise}

\begin{exercise}[$\star\star$]\label{ex:two}
\emph{How $2$ splits.} Show that $(2)=P^2$ is ramified when $d\equiv2,3\pmod4$, with $P=(2,\sqrt d)$ or $P=(2,1+\sqrt d)$; that $(2)=P\bar P$ with $P\ne\bar P$ when $d\equiv1\pmod8$, with $P=(2,\omega)$; and that $(2)$ is maximal when $d\equiv5\pmod8$.
\end{exercise}

\begin{exercise}[$\star\star$]\label{ex:factor5}
In $\Z[\sqrt{-5}]$, factor the ideals $(2)$, $(3)$, $(5)$, $(7)$, $(11)$ and $(29)$ into maximal ideals, and say which factors are principal.
\end{exercise}

\begin{exercise}[$\star\star$]\label{ex:ufdh1}
Conversely to Proposition~\ref{prop:h1}, show that if $\mathcal O$ has unique factorisation into irreducibles, then every ideal is principal, so $h_K=1$.
\end{exercise}

\begin{exercise}[$\star\star\star$]\label{ex:x2p5}
Show that $x^2+5=y^3$ has no solution in integers. \emph{Hint: factor $(x+\sqrt{-5})(x-\sqrt{-5})$ into ideals and use $h=2$.}
\end{exercise}

\begin{exercise}[$\star\star\star$]\label{ex:fermat5}
\emph{Euler's conjecture, proved by Lagrange.} Show that a prime $p\ne5$ is of the form $x^2+5y^2$ with $x,y\in\Z$ if and only if $p\equiv1$ or $9\pmod{20}$.
\end{exercise}

\begin{beauty}
Theorem~\ref{thm:dedekind} is the moment where number theory becomes algebra. The failure of unique factorisation, which seemed fatal, is measured by a finite abelian group, and in examples like Exercises~\ref{ex:x2p5} and~\ref{ex:fermat5} one can still argue as if factorisation were unique, as long as one tracks the class group. The pattern of Theorem~\ref{thm:split}, where the splitting of $p$ depends only on $p$ modulo $4|d|$, is quadratic reciprocity seen from inside $\mathcal O$; its generalisation to all abelian extensions of $\Q$ and beyond is class field theory. In Exercise~\ref{ex:fermat5} the answer is still a congruence ($p\equiv1,9\pmod{20}$), but the proof already needs the class group, since $x^2+5y^2$ and $2x^2+2xy+3y^2$ share the same discriminant. For some $n$, such as $n=14$ or $n=27$, no congruence condition on $p$ describes the primes $x^2+ny^2$ at all: this is the subject of Cox's beautiful book \emph{Primes of the form $x^2+ny^2$}.
\end{beauty}

\section{Hints}

\begin{hint}{ex:sqrt10}
Reduce $a^2-10b^2$ modulo $5$.
\end{hint}

\begin{hint}{ex:two}
Lemma~\ref{lem:hurwitz} does all the computations. In the last case, compute $N(m+k\omega)$ modulo $2$.
\end{hint}

\begin{hint}{ex:factor5}
Look for square roots of $-5$ modulo $p$. A factor of norm $p$ is principal exactly when $p$ is a norm $x^2+5y^2$.
\end{hint}

\begin{hint}{ex:ufdh1}
A maximal ideal contains an irreducible $\pi$. To see that $(\pi)$ is maximal, use that $\mathcal O/(\pi)$ is finite, so the powers of an element must repeat.
\end{hint}

\begin{hint}{ex:x2p5}
First prove $\gcd(y,10)=1$ with congruences modulo $8$ and $25$.
\end{hint}

\begin{hint}{ex:fermat5}
If $P$ is not principal, then $PP_2$ is, where $P_2=(2,1+\sqrt{-5})$. What is its norm?
\end{hint}

\section{Solutions}

\begin{solution}{ex:sqrt10}
\idea Irreducibility: a proper factor of an element of norm $\pm4$, $\pm25$ or $-10$ must have norm $\pm2$ or $\pm5$, and the form $a^2-10b^2$ never takes these values, which a reduction modulo $5$ shows. Then two different factorisations contradict Proposition~\ref{prop:h1}.
\writeup $10\equiv2\pmod4$, so $\mathcal O=\Z[\sqrt{10}]$, with $N(a+b\sqrt{10})=a^2-10b^2$. If $a^2-10b^2=\pm2$, then $a^2\equiv\pm2\pmod5$, impossible (squares are $0,\pm1$). If $a^2-10b^2=\pm5$, then $5\mid a$, and writing $a=5a'$, $5a'^2-2b^2=\pm1$, so $-2b^2\equiv\pm1$, $b^2\equiv\mp3\equiv2$ or $3\pmod5$, impossible. Now $N2=4$, $N5=25$, $N\sqrt{10}=-10$, so a factorisation into two non-units would contain a factor of norm $\pm2$ or $\pm5$: all three are irreducible. They are pairwise non-associate since their norms have different absolute values. So $10=2\cdot5=\sqrt{10}\cdot\sqrt{10}$ are two different factorisations, and $h\ge2$ by Proposition~\ref{prop:h1}. (In ideals: $(2)=P^2$, $(5)=Q^2$, $(\sqrt{10})=PQ$ with $P=(2,\sqrt{10})$, $Q=(5,\sqrt{10})$, by Exercise~\ref{ex:two} and Theorem~\ref{thm:split}(b).)
\end{solution}

\begin{solution}{ex:two}
\idea Imitate the proof of Theorem~\ref{thm:split}: compute $P\bar P$ with Lemma~\ref{lem:hurwitz}, and in the inert case look at the norm of an element modulo $2$.
\writeup If $d\equiv2\pmod4$, $P=(2,\sqrt d)=\bar P$ and Lemma~\ref{lem:hurwitz} gives $P^2=(\gcd(4,-d,0))=(2)$. If $d\equiv3\pmod4$, $P=(2,1+\sqrt d)$ equals $\bar P=(2,1-\sqrt d)$ since $1-\sqrt d=2-(1+\sqrt d)$, and $P^2=(\gcd(4,1-d,\operatorname{Tr}(2(1-\sqrt d))))=(\gcd(4,1-d,4))=(2)$ because $1-d\equiv2\pmod4$. If $d\equiv1\pmod8$, $N\omega=\frac{1-d}4$ is even, and $P=(2,\omega)$ has $P\bar P=(\gcd(4,\frac{1-d}4,\operatorname{Tr}(2\bar\omega)))=(\gcd(4,\tfrac{1-d}4,2))=(2)$; and $P\ne\bar P$, otherwise $P\ni\omega+\bar\omega=1$. If $d\equiv5\pmod8$, $\frac{1-d}4$ is odd and $N(m+k\omega)=m^2+mk+\frac{1-d}4k^2\equiv m^2+mk+k^2\pmod2$, which is odd unless $m,k$ are both even. As in Theorem~\ref{thm:split}(c), an ideal $J\supseteq(2)$ of norm $2$ would contain some $\alpha=m+k\omega\notin2\mathcal O$ with $N\alpha$ even, which we just excluded; so $(2)$ is maximal. In the ramified and split cases the displayed maximal factors have norm $2$. In the inert case the single maximal ideal $(2)$ has norm $4$. Maximality of norm-$2$ factors follows from Lemma~\ref{lem:maximal}(c); inert maximality was proved separately.
\end{solution}

\begin{solution}{ex:factor5}
\idea Use Theorem~\ref{thm:split}: compute $\left(\frac{-5}p\right)$, find a square root $a$ of $-5$ modulo $p$, and write $P=(p,a+\sqrt{-5})$. A factor $P$ of norm $p$ is principal if and only if $p=N\alpha=x^2+5y^2$ for some $\alpha$, because $N((\alpha))=N\alpha$.
\writeup First the criterion. If $P=(\alpha)$ then $N\alpha=N(P)=p$. Conversely, if $N\alpha=p$, then $(\alpha)$ has norm $p$, contains $p=\alpha\bar\alpha$, so divides $(p)=P\bar P$, and is $P$ or $\bar P$; its conjugate $(\bar\alpha)$ is the other one, so both are principal.

$(2)=P^2$ with $P=(2,1+\sqrt{-5})$, not principal (Proposition~\ref{prop:six}). $(3)=Q\bar Q$ with $Q=(3,1+\sqrt{-5})$ ($1\equiv-5\pmod3$), not principal since $x^2+5y^2\ne3$. $(5)=(5,\sqrt{-5})^2=(\sqrt{-5})^2$, principal, as $5=-\sqrt{-5}^2$. For $7$: $-5\equiv2\equiv3^2\pmod7$, so $(7)=R\bar R$ with $R=(7,3+\sqrt{-5})$, not principal ($x^2+5y^2=7$ has no solution: $y=0$ or $y=\pm1$ give $7$ or $2$ as a square). For $11$: $-5\equiv6\pmod{11}$, and the squares modulo $11$ are $1,4,9,5,3$, so $\left(\frac{-5}{11}\right)=-1$ and $(11)$ is maximal (and principal). For $29$: $13^2=169=5\cdot29+24\equiv-5$, so $(29)=S\bar S$ with $S=(29,13+\sqrt{-5})$, and $29=3^2+5\cdot2^2$ shows both factors are principal. Which is which? $3+2\sqrt{-5}=29-2(13-\sqrt{-5})\in\bar S$, so $\bar S=(3+2\sqrt{-5})$ and $S=(3-2\sqrt{-5})$.
\end{solution}

\begin{solution}{ex:ufdh1}
\idea Since every ideal is a product of maximal ideals, it suffices to show each maximal $P$ is principal. $P$ contains a non-zero integer, hence one of its irreducible factors $\pi$. Unique factorisation makes $\pi$ prime, and then $(\pi)$ is maximal, by a finiteness argument on the powers of an element modulo $\pi$. So $P=(\pi)$.
\writeup First, an irreducible $\pi$ is prime: if $\pi\mid ab$, write $ab=\pi c$ and factor $a,b,c$; by uniqueness $\pi$ is an associate of an irreducible factor of $a$ or of $b$. Let $P$ be maximal. It contains $N(P)\ge2$, which factors into irreducibles, so by Lemma~\ref{lem:maximal}(a) $P$ contains an irreducible $\pi$, and $(\pi)\subseteq P$. We show $(\pi)$ is maximal. Let $J\supsetneq(\pi)$ and $a\in J\setminus(\pi)$. The ideal $(\pi)$ contains $N\pi\ne0$, so $\mathcal O/(\pi)$ is finite (Proposition~\ref{prop:idealrules}(d)), and $a^i\equiv a^j\pmod\pi$ for some $i<j$. Then $\pi\mid a^i(a^{j-i}-1)$, and since $\pi$ is prime and $\pi\nmid a$, $\pi\mid a^{j-i}-1$. So $1=a^{j-i}-\pi z\in J$, and $J=\mathcal O$. Hence $P=(\pi)$. By Theorem~\ref{thm:dedekind} every ideal is a product of principal ideals, hence principal, and $h_K=1$.
\end{solution}

\begin{solution}{ex:x2p5}
\idea Follow Exercise~\ref{ex:x3y21}, with ideals. Congruences give $\gcd(y,10)=1$; then the ideals $(x\pm\sqrt{-5})$ have no common maximal factor, so each is the cube of an ideal $I$. The class of $I$ has order dividing $3$ and dividing $2$ (the class group has order $2$), so $I$ is principal, and we are back to comparing coefficients of a cube.
\writeup If $x$ is odd, $x^2+5\equiv6\pmod8$, which is not a cube modulo $8$ (an even cube is $\equiv0$, an odd one is odd). So $x$ is even and $y$ odd. If $5\mid y$, then $5\mid x$, and $x^2+5\equiv5\pmod{25}$ while $y^3\equiv0$: impossible. So $\gcd(y,10)=1$.

Let $A=(x+\sqrt{-5})$, $\bar A=(x-\sqrt{-5})$, so $A\bar A=(y)^3$. A maximal $P$ containing $A$ and $\bar A$ contains $2\sqrt{-5}$, hence $(2\sqrt{-5})^2=-20$, and contains $y^3$; since $\gcd(20,y^3)=1$ it contains $1$: impossible. Write $(y)=\prod P_i^{e_i}$. Each $P_i$ contains $A\bar A$, hence exactly one of $A$, $\bar A$ (Lemma~\ref{lem:maximal}(a)), and by uniqueness in Theorem~\ref{thm:dedekind} all of $P_i^{3e_i}$ goes into that one. So $A=I^3$ for an ideal $I$. In the class group of order $2$ (Theorem~\ref{thm:hm5}), $[I]^2=1$ and $[I]^3=[A]=1$, so $[I]=1$: $I=(\gamma)$, and $x+\sqrt{-5}=u\gamma^3$ with $u=\pm1$, a cube. So $x+\sqrt{-5}=(a+b\sqrt{-5})^3=(a^3-15ab^2)+(3a^2b-5b^3)\sqrt{-5}$, and $b(3a^2-5b^2)=1$. Then $b=\pm1$ and $3a^2-5=\pm1$, so $a^2\in\{2,\frac43\}$: no solution.
\end{solution}

\begin{solution}{ex:fermat5}
\idea One direction is congruences. For the other, $p\equiv1,9\pmod{20}$ makes $-5$ a square modulo $p$, so $(p)=P\bar P$ with $N(P)=p$. By Theorem~\ref{thm:hm5}, $P$ is principal (done, $p=x^2+5y^2$) or equivalent to $P_2=(2,1+\sqrt{-5})$. In the second case $PP_2$ is principal of norm $2p$, so $2p=a^2+5b^2$, and a congruence modulo $5$ rules this out.
\writeup If $p=x^2+5y^2$ with $p\ne5$, then $p\ne2$ ($2$ is not of this form), so $x$ and $y$ have opposite parity, $p\equiv x^2+y^2\equiv1\pmod4$, and $p\equiv x^2\equiv\pm1\pmod5$ since $5\nmid x$. Hence $p\equiv1$ or $9\pmod{20}$.

Conversely let $p\equiv1,9\pmod{20}$. Then $\left(\frac{-5}p\right)=\left(\frac{-1}p\right)\left(\frac5p\right)=\left(\frac p5\right)=1$, by Euler's criterion (Theorem~\ref{thm:eulercrit}), quadratic reciprocity (Theorem~\ref{thm:qr}) and $p\equiv\pm1\pmod5$. By Theorem~\ref{thm:split}, $(p)=P\bar P$ with $N(P)=p$. If $P$ were not principal, it would be equivalent to $P_2$ (Theorem~\ref{thm:hm5}), so $PP_2\sim P_2^2=(2)$ would be principal, $PP_2=(\alpha)$ with $N\alpha=N(P)N(P_2)=2p$. Then $a^2+5b^2=2p$ gives $2p\equiv a^2\equiv\pm1\pmod5$ (as $5\nmid 2p$), i.e.\ $p\equiv\pm2\pmod5$, a contradiction. So $P=(x+y\sqrt{-5})$ and $p=N(P)=x^2+5y^2$. For example $29=3^2+5\cdot2^2$, $41=6^2+5\cdot1^2$, $61=4^2+5\cdot3^2$, while $3,7,23,43$ (which split, but into non-principal ideals) satisfy $2p=a^2+5b^2$ instead: $6=1+5$, $14=9+5$, $46=1+45$, $86=81+5$.
\end{solution}

\section*{Sources and further reading}

\begin{itemize}
\item R.~Dedekind, \emph{Theory of Algebraic Integers} (1877), translated by J.~Stillwell, Cambridge University Press, 1996. The original source of ideals, still a pleasure to read.
\item H.~M.~Edwards, \emph{Fermat's Last Theorem: A Genetic Introduction to Algebraic Number Theory}, Graduate Texts in Mathematics~50, Springer, 1977. Kummer's ideal numbers, told historically.
\item K.~Ireland and M.~Rosen, \emph{A Classical Introduction to Modern Number Theory}, 2nd ed., Graduate Texts in Mathematics~84, Springer, 1990, Chapters~12--13. Rings of integers and the splitting of primes.
\item P.~Samuel, \emph{Th\'eorie alg\'ebrique des nombres}, Hermann, 1967 (English translation: \emph{Algebraic Theory of Numbers}, Dover, 2008). Short and clear; the natural next step after this chapter.
\item I.~Stewart and D.~Tall, \emph{Algebraic Number Theory and Fermat's Last Theorem}, 4th ed., CRC Press, 2016. Class groups, Kummer's theorem on regular primes, and the story up to Wiles.
\item D.~A.~Marcus, \emph{Number Fields}, 2nd ed., Universitext, Springer, 2018. The general theory, with many exercises.
\item D.~A.~Cox, \emph{Primes of the Form $x^2+ny^2$}, 3rd ed., AMS Chelsea, 2022. Exercise~\ref{ex:fermat5} and its generalisations, and the history of the class number~$1$ problem.
\item H.~M.~Edwards, \emph{Fermat's Last Theorem}, Springer, 1977, Chapter~4. Lam\'e's announcement of 1 March 1847, Liouville's objection, and Kummer's letter.
\end{itemize}
